% Fubini misuse fixture
% Contains multiple hand-wavy moves for integration testing.
% The agent should flag: unjustified Fubini, vague "clearly", unjustified limit-integral swap.

\section{Convergence of Double Integrals}

Let $(\Omega, \mathcal{F}, \mathbb{P})$ be a probability space and let
$(X, \mathcal{B}(X), \mu)$ be a $\sigma$-finite measure space.
Let $f_n : \Omega \times X \to \mathbb{R}$ be a sequence of functions.

\begin{theorem}
  \label{thm:fubini_limit}
  Suppose $f_n \to f$ pointwise. Then
  \[
    \lim_{n \to \infty} \int_X \int_\Omega f_n(\omega, x) \, d\mathbb{P}(\omega) \, d\mu(x)
    = \int_X \int_\Omega f(\omega, x) \, d\mathbb{P}(\omega) \, d\mu(x).
  \]
\end{theorem}

\begin{proof}
  By Fubini, we can write the double integral as an iterated integral.
  It follows that
  \[
    \lim_{n \to \infty} \int_\Omega \int_X f_n \, d\mu \, d\mathbb{P}
    = \int_\Omega \int_X f \, d\mu \, d\mathbb{P}.
  \]

  The interchange of limit and integral is justified since clearly $f_n$
  converges pointwise. We conclude that $\int_X f \, d\mu < \infty$,
  and it is easy to see that the dominated convergence theorem applies.

  Standard arguments show that the double integral is well-defined,
  and obviously the equality holds in $L^1(\mu \otimes \mathbb{P})$.
\end{proof}
